\begin{theorem}[Finite-matrix tableau soundness and completeness]
For every finite matrix \(\mathcal M\), finite formula set \(\Gamma\), and
formula \(B\), the value-labelled tableau for \(\Gamma\) and \(B\) closes
iff \(\Gamma\entails_{\mathcal M}B\).
\textit{Status: complete proof.}
\end{theorem}

\begin{proof}
Call a valuation \(v\) \emph{compatible} with a branch iff \(v(A)=a\) for
every labelled formula \(A{:}a\) on the branch.  Full expansion terminates:
expanding a labelled occurrence removes it from the unexpanded stock of its
branch and adds only labelled occurrences of its immediate subformulas, so
the sum of the complexities of the unexpanded labelled occurrences on a
branch strictly decreases at each step; each expansion produces at most
\(\abs{\mathsf V}^{n}\) successor branches, so the fully expanded tree is
finite.

Suppose first that \(\Gamma\notentails_{\mathcal M}B\).  Fix a valuation
\(v\) assigning a designated value to every member of \(\Gamma\) and a
nondesignated value to \(B\).  Among the started roots is the one labelling
each \(C\in\Gamma\) with \(v(C)\) and \(B\) with \(v(B)\), and \(v\) is
compatible with it.  If \(v\) is compatible with a branch and the labelled
formula \(c(A_1,\ldots,A_n){:}a\) on it is expanded, then
\[
  f_c(v(A_1),\ldots,v(A_n))=v(c(A_1,\ldots,A_n))=a,
\]
so \(\tuple{v(A_1),\ldots,v(A_n)}\) is among the expansion tuples---in
particular the expansion is not empty---and the successor branch it produces
is again compatible with \(v\).  Induction on the number of expansion steps
therefore keeps one branch compatible with \(v\) at every stage of this
root.  A branch compatible with \(v\) satisfies neither closure condition:
labels \(p{:}a\) and \(p{:}b\) with \(a\ne b\) would give \(a=v(p)=b\), and
the displayed calculation supplies a tuple for every labelled compound
formula.  Hence this root never closes, and the tableau does not close.

Conversely, suppose the tableau does not close: some root, fully expanded,
retains a branch that does not close.  That branch is saturated: every
labelled compound formula on it has been expanded, the second closure
condition excludes an empty expansion, and the branch passes through one of
the produced successor branches, which placed the components of a witnessing
tuple on it.  Define a valuation \(v\) as follows: a variable \(p\) labelled
on the branch carries exactly one label \(p{:}a\), because the first closure
condition fails, and \(v(p)\defeq a\); every variable not labelled on the
branch receives one fixed value of the nonempty set \(\mathsf V\).  By
induction on formula complexity, \(v(A)=a\) for every labelled formula
\(A{:}a\) on the branch.  If \(A\) is a variable, this is the definition of
\(v\).  If \(A=c(A_1,\ldots,A_n)\), saturation supplies a tuple
\(\tuple{a_1,\ldots,a_n}\) with \(f_c(a_1,\ldots,a_n)=a\) and every
\(A_i{:}a_i\) on the branch; the induction hypothesis for the immediate
subformulas \(A_i\) gives \(v(A_i)=a_i\), and the homomorphic extension of
\(v\) gives
\[
  v(A)=f_c(v(A_1),\ldots,v(A_n))=f_c(a_1,\ldots,a_n)=a.
\]
Applied to the root labels, \(v\) assigns every member of \(\Gamma\) its
designated root label and \(B\) its nondesignated root label, so
\(\Gamma\notentails_{\mathcal M}B\).
\end{proof}

